
Sedmý případ užití je specifikován v tabulce č. \ref{tab:usecase:UC07}.

\begin{UseCaseTable}
  {Změna parametrů s časovou platností}
  {UC07}
  {F07}
  {Aktualizace parametrů subjektu zájmu nebo senzoru s platností od zadaného data při zachování návaznosti starších měření na původní parametry}
  {Provozovatel monitorovaného prostoru}
  {}
  {Uživatel je přihlášen, subjekt zájmu (budova, místnost) nebo senzor je v systému evidován}
  {V evidenci je uložena nová verze parametrů s platností od zadaného data, předchozí verze parametrů zůstává platná pro období do tohoto data, historická měření zůstávají vázána na verzi parametrů platnou v době jejich pořízení}

\UseCaseScenario{Základní scénář}
\UseCaseStep{uživatel}{V detailu evidovaného subjektu zájmu nebo senzoru zvolí akci pro úpravu parametrů.}
\UseCaseStep{systém}{Vrátí formulář s aktuálně platnými parametry a polem pro datum, od kdy mají nové hodnoty platit.}
\UseCaseStep{uživatel}{Upraví hodnoty parametrů, zadá datum platnosti nových hodnot a formulář odešle.}
\UseCaseStep{systém}{Ověří, že datum platnosti je zadáno a všechny zadané hodnoty jsou v platném tvaru.}
\UseCaseStep{systém}{Uloží novou verzi parametrů jako platnou od zadaného data, předchozí verzi zachová pro období do tohoto data a zobrazí potvrzení.}

\UseCaseScenario{Alternativní scénář -- neplatné údaje}
\UseCaseStepCustom{A1}{uživatel, systém}{shodné s kroky 1--3 základního scénáře.}
\UseCaseStepCustom{A2}{systém}{Zjistí chybu ve vyplnění formuláře (chybí datum platnosti, neplatný formát hodnoty) a vyzve uživatele k její opravě.}
\UseCaseStepCustom{A3}{uživatel}{Opraví zadané údaje.}
\UseCaseStepCustom{A4}{systém}{Pokračuje krokem 4 základního scénáře.}

\end{UseCaseTable}
